\section{Generation States and a Closure-Defect Coupling}
\label{v480:generation-coupling}

This section constructs a conditional operator interface from objects
already named in the paper: the generation slots, the closure projector
and the shifted pacer. It provides explicit formulas, not yet evaluated
native matrix elements or a derived Yukawa interaction. In particular,
it separates a mass readout in fixed slots from one in coupling
eigenstates. Neither is identified with physical lepton masses merely
by giving it that name.

\subsection{Closure-Compatible Generation States}
\label{v480:generation-states}

Use the positive-metric real Hilbert realization from
Section~\ref{v480:hc-mass-identification}. For a nonzero odd sector,
one may choose a unit vector $\phi_-=P_-g/\|P_-g\|$ with
$P_-g\ne0$. The inherited Dixon dictionary
(Theorem~\ref{thm:three-generations}) suggests the normalized seed
slots $\phi_-\otimes e_4$, $\phi_-\otimes e_5$,
$\phi_-\otimes e_6$. This requires a declared tensor-product metric
and a generation-neutral extension of the real $J$-action; the
dictionary alone does not prove that factorization in a native carrier.

On the complexification of this real realization, let
$T:\mathbb C^3\to\mathcal K$ isometrically embed those slots.
Here $J$ denotes the complex-linear extension of the real involution,
not an unannounced identification with the inherited anti-unitary
convention. Realize $P=\hat P_{C_7}$ as an orthogonal closure
projector on $\mathcal K$, commuting with this $J$ on the declared
realization. It is \emph{not} the HC projector $C$. An action of the
pacer on this same carrier must still be supplied; a bosonic clock
does not automatically define its action on the probe slots.

\begin{proposition}[Normalized closed generation seeds]
\label{v480:generation-closed-seeds}
If $G=T^*PT$ is positive definite, then
\[
 S=PTG^{-1/2},\qquad S^*S=I_3,\quad PS=S,\quad JS=-S.
\]
Thus $\psi_i=S\varepsilon_i$ are three orthonormal closed odd states,
where $\varepsilon_i$ is the standard generation-coordinate basis.
If $G$ has rank less than three, this construction cannot produce
three independent closed states from these seeds.
\end{proposition}
\begin{proof}
The positive inverse square root exists in the three-dimensional
seed space. Since $P^2=P=P^*$,
$S^*S=G^{-1/2}T^*PTG^{-1/2}=I_3$ and $PS=S$.
Commutation with $J$ and $JT=-T$ give $JS=-S$.
Multiplication by $P$ cannot restore a dimension lost in $PT$.
\end{proof}

The choice of $g$, the common carrier and $P$ remain native input
obligations. If $G$ is not diagonal in the seed slots, this
orthonormalization already mixes them; it is not a proof that the
original octonion labels are physical mass eigenstates.

\paragraph{The spectral restriction is not the new coupling.}
The inherited proposal $D_-=L_0|_{H_-}$ must not be identified with
the new mass coupling $D_\beta$ by its letter. On the real grading
used here, a sector-preserving real-linear $A:H_-\to H_-$ commutes
with $J=-I$ there. Requiring also $AJ=-JA$ forces $A=0$.
For the inclusion $\iota:H_-\hookrightarrow H$, the correctly typed
relation $J\iota A=-\iota A$ only records that the output is
odd, not anticommutation. This qualifies the D1 argument in
Theorem~\ref{thm:op2-existence}; its nonzero restriction cannot be
borrowed as an already verified anticommuting mass coupling.
Domain and self-adjointness questions for that spectral proposal
are not settled by the finite construction here.

\subsection{Constructing the Coupling From the Pacer and Closure}
\label{v480:generation-defect}

Fix a bounded realization of $R_\beta$ on the common carrier, or a
declared finite matrix realization. In the inherited notation it is
the proposed $\hat{\mathcal R}_\beta$ of
Theorem~\ref{thm:dirac-probe}. Fix a positive reference $r_{\rm ref}$
with the same units as $R_\beta$, and put
$\widetilde R_\beta=R_\beta/r_{\rm ref}$. These domain/unit data
are additional specifications, not inferred from a clock label.

\begin{proposition}[A positive closure-defect readout]
\label{v480:generation-defect-gram}
For the closed embedding above, define
\[
 D_\beta=(I-P)\widetilde R_\beta P,\qquad
 B_\beta=D_\beta S,\qquad Q_\beta=B_\beta^*B_\beta.
\]
Then $D_\beta S=[\widetilde R_\beta,P]S$ and
$Q_\beta$ is positive semidefinite. Its entries are
\[
 (Q_\beta)_{ij}
 =\langle (I-P)\widetilde R_\beta\psi_i,
                 (I-P)\widetilde R_\beta\psi_j\rangle.
\]
For every $P\psi=\psi$,
$\langle\psi,[\widetilde R_\beta,P]\psi\rangle=0$,
even though $D_\beta\psi$ may be nonzero.
\end{proposition}
\begin{proof}
On a closed input,
$[\widetilde R_\beta,P]\psi=(I-P)\widetilde R_\beta\psi$,
which is orthogonal to $\psi\in\operatorname{im}P$.
For $v\in\mathbb C^3$,
$\langle v,Q_\beta v\rangle=\|B_\beta v\|^2\geq0$,
and the matrix-entry formula is the defining Gram identity.
\end{proof}

Thus the commutator \emph{expectation} used for mass in
Proposition~\ref{prop:three-compatible-readings} vanishes on an
exactly closed state under these projector assumptions. Replacing it
by a transition norm is a genuine change of readout, not an equivalent
rewriting of that inherited claim. The norm measures a tendency to
leave the selected closure space under the pacer; it does not prove
that a physical state actually leaves it or that such leakage is a
derived mass mechanism. In particular, exact $[R_\beta,P]=0$ makes
this candidate vanish. A common mass scale $\kappa>0$ is still needed.

\subsection{Coordinate Readouts and Candidate Mass Eigenstates}
\label{v480:generation-spectrum}

\begin{proposition}[The Gram spectrum and the slot diagonal differ]
\label{v480:generation-mass-spectrum}
Let $Q_\beta=V\operatorname{diag}(q_0,q_1,q_2)V^*$,
with $V$ unitary and $q_i\geq0$. The states
$\chi_i=SV\varepsilon_i$ are orthonormal closed odd states and
$\|D_\beta\chi_i\|^2=q_i$. If one additionally postulates the
effective mass matrix $M_\beta=\kappa Q_\beta$, its candidate
eigenmasses are $\kappa q_i$, whereas the fixed-slot norm readouts
are $\kappa(Q_\beta)_{ii}$. When $Q_\beta\ne0$, the spectral
Koide diagnostic is
\[
 K_{\rm spec}
 =\frac{\operatorname{Tr}Q_\beta}
        {(\operatorname{Tr}\sqrt{Q_\beta})^2}.
\]
It is invariant under unitary changes of generation coordinates.
\end{proposition}
\begin{proof}
Use $S^*S=I_3$ and $V^*Q_\beta V=\operatorname{diag}(q_i)$.
The principal root masses are $\sqrt\kappa\sqrt{q_i}$; the
common scale cancels. Trace and functional calculus are unitarily
invariant. The fixed-slot formula is the diagonal Gram identity.
\end{proof}

For example, $Q=\tfrac13\mathbf1\mathbf1^*$, with
$\mathbf1=(1,1,1)^{\mathsf T}$, has three equal diagonal entries
but eigenvalues $(1,0,0)$. Its slot diagnostic is $1/3$ and its
spectral diagnostic is $1$. Equal diagonal readouts do not establish
degenerate eigenmasses. Conversely, diagonalizing a positive matrix
does not identify it with the native first-order mass operator
$\Omega$ of Theorem~\ref{thm:master-field-equation}; a physical
derivation must justify $M_\beta=\kappa Q_\beta$ and its units,
rather than confuse that hypothesis with $\Omega^*\Omega$.

\subsection{A Cyclic Coupling That Produces the Brannen Form}
\label{v480:generation-cyclic}

There is a precise sufficient link to a generation action on the
carrier. If a unitary $W$ satisfies $WS=SU$, $WP=PW$ and
$W\widetilde R_\beta=\widetilde R_\beta W$, then
$B_\beta U=WB_\beta$, whence $Q_\beta U=UQ_\beta$.
These follow by moving $W$ through the defining projectors and
pacer and then taking adjoints. The three compatibility relations
must be supplied for the native realization, not inferred merely
from the existence of three seed slots.

\begin{proposition}[A Hermitian cyclic root coupling]
\label{v480:generation-cyclic-root}
Let $U\varepsilon_i=\varepsilon_{i+1\bmod3}$ and set
\[
 L=aI_3+zU+\bar zU^*,\qquad a>0,\quad z=r e^{i\theta},\ r\geq0.
\]
Choose the Fourier labels so $U$ has eigenvalues
$e^{2\pi i k/3}$. Then $L$ has eigenvalues
\[
 \ell_k=a+2r\cos(\theta+2\pi k/3).
\]
If all $\ell_k\geq0$, and if the closure-defect Gram matrix is
independently established to be $Q=L^2$, its candidate masses are
$m_k=\kappa\ell_k^2$ and
\[
 K_{\rm spec}=\frac13+\frac{2r^2}{3a^2}.
 \qquad K_{\rm spec}=\frac23\ \Longleftrightarrow\ r=a/\sqrt2.
\]
Conversely, every nonzero positive semidefinite $Q$ commuting with
$U$ has a positive root $L=\sqrt Q$ of the displayed form, with
\[
 a=\tfrac13\operatorname{Tr}\sqrt Q,\qquad
 z=\tfrac13\operatorname{Tr}(U^*\sqrt Q).
\]
\end{proposition}
\begin{proof}
Fourier diagonalization gives the eigenvalues. Since
$\operatorname{Tr}U=\operatorname{Tr}U^2=0$,
$\operatorname{Tr}L=3a$ and
$\operatorname{Tr}L^2=3a^2+6r^2$.
Positivity gives $\sqrt{L^2}=L$, so division gives the result.
Conversely, $\sqrt Q$ commutes with $U$ by spectral calculus.
The three distinct Fourier eigenvalues of $U$ make its commutant
the span of $I_3,U,U^2$. Hermiticity and $U^*=U^2$ give the stated
coefficient form; the trace identities recover $a$ and $z$.
\end{proof}

This constructs the familiar three-phase harmonic form as eigenvalues
of a concrete coupling matrix, without identifying three raw triolic
coordinates with mass roots. It does not derive $r/a=1/\sqrt2$,
$\theta$ or $\kappa$, nor does it establish cyclic compatibility
or the coefficients of the native $B_\beta^*B_\beta$. Once those
native data are given, the converse specifies the coefficients
without a mass fit. For negative $\ell_k$ in a proposed $L$,
the root trace is $\sum_k|\ell_k|$ instead, and the displayed Koide
formula is not valid. For real $z$, two $\ell_k$ coincide; a generic
complex relative coefficient can lift this pair degeneracy. That is
an internal generation phase, not an overall phase of the pacer.

There is no conflict with Proposition~\ref{v480:hc-mass-equivariant-obstruction}:
the seed slots form a cyclic \emph{orbit}, not one invariant input,
and the spectral states are Fourier combinations of those slots.
A complex Hermitian circulant matrix may have three different
eigenvalues while its diagonal entries are all equal. If one insists
on a real symmetric cyclic matrix, it is of the form
$a_0I+b_0(U+U^*)$ and has a doubly degenerate eigenvalue.
The choice of scalar field and physical generation basis matters.

\begin{proposition}[Positive phase chambers and a small-root boundary]
\label{v480:generation-phase-chambers}
At $r=a/\sqrt2$, the cyclic root matrix is positive semidefinite
exactly for the phase chambers
\[
 \theta\in\bigcup_{j=0}^2
 \left[\frac{2\pi j}{3}-\frac{\pi}{12},
       \frac{2\pi j}{3}+\frac{\pi}{12}\right]
 \pmod{2\pi}.
\]
Near the upper boundary $\theta_b=3\pi/4$ of the chamber centered
at $2\pi/3$, put $\theta=\theta_b-\delta$, $\delta\geq0$.
Then $\ell_0=a\delta+O(\delta^2)$ and
$m_0=\kappa a^2\delta^2+O(\delta^3)$, while the other two
root eigenvalues tend to $a(3-\sqrt3)/2$ and $a(3+\sqrt3)/2$.
\end{proposition}
\begin{proof}
Each $\ell_k\geq0$ requires
$\cos(\theta+2\pi k/3)\geq-1/\sqrt2$.
Intersecting these three conditions in each cyclic period gives
the displayed intervals. At $3\pi/4$ the zeroth root vanishes
and its phase derivative is $-a$. Taylor expansion and squaring
give the small-mass law. Substitution gives the other limits.
\end{proof}

A relative generation phase near such a boundary can therefore
produce a parametrically light eigenmass within this candidate.
This explains a possible amplification mechanism, not why a native
state should select that phase or the Koide radius in the first place.

\subsection{Overall Phase, Relative Phase and Underdetermination}
\label{v480:generation-phase}

\begin{proposition}[A common pacer phase cancels]
\label{v480:generation-global-phase}
For fixed $P$, $S$ and normalization, if
$R_\beta=e^{i\varphi(\beta)}R_0$, then
$B_\beta=e^{i\varphi(\beta)}B_0$ and $Q_\beta=Q_0$.
Both slot and spectral norm-based mass readouts are unchanged.
\end{proposition}
\begin{proof}
Pull the scalar through the projectors and use its unit modulus in
$B_\beta^*B_\beta$.
\end{proof}

The written pacer $\hat R_{\rm GUE}e^{i\pi\bsym}$ in
Theorem~\ref{thm:dirac-probe} has exactly this common-phase form.
It supplies no $\beta$-dependent mass ratios for this norm readout
when $P$ and $S$ are fixed. Allowing them to vary requires their
own independently specified frame laws.
The absolute commutator expectation in the inherited mass formula
also loses this common phase, even on inputs that are not exactly
closed: its expectation acquires only the same unit-modulus factor.

\begin{proposition}[A two-channel relative phase can survive]
\label{v480:generation-relative-phase}
As an additional ansatz, let
$\widetilde R_\beta=R_a+e^{i\varphi(\beta)}R_b$ with matched
domains and units. Put $A=(I-P)R_aS$, $B=(I-P)R_bS$. Then
\[
 Q_\beta=A^*A+B^*B
       +e^{i\varphi(\beta)}A^*B
       +e^{-i\varphi(\beta)}B^*A.
\]
The phase cancels when $A^*B=0$; it can otherwise affect the Gram
matrix through interference. A varying matrix alone does not imply
varying eigenmass ratios: the variation could be a basis conjugation.
\end{proposition}
\begin{proof}
Expand $(A+e^{i\varphi}B)^*(A+e^{i\varphi}B)$.
Unitary conjugation preserves its eigenvalues.
\end{proof}

The triadic channels in Proposition~\ref{prop:triadic-two-probe}
are a place to seek such extra matrix elements, not a specification
of them. They must be realized on the same generation/probe carrier
and converted to compatible units before adding them to a clock
operator. A term of the form $PXP$ cannot itself contribute to
$(I-P)\widetilde R_\beta P$. The internal coupling must actually
connect the selected and complementary spaces for this candidate
to be nonzero.

\begin{proposition}[Even unitary pacers leave the Gram matrix free]
\label{v480:generation-unitary-dilation}
Every three-dimensional positive contraction $0\leq Q\leq I_3$
can occur as the closed-to-complementary Gram matrix of a unitary
pacer. On $\mathbb C^3\oplus\mathbb C^3$, take
\[
 P=\begin{pmatrix}I_3&0\\0&0\end{pmatrix},\quad
 S=\begin{pmatrix}I_3\\0\end{pmatrix},\quad
 R=\begin{pmatrix}\sqrt{I_3-Q}&-\sqrt Q\\
                   \sqrt Q&\sqrt{I_3-Q}\end{pmatrix}.
\]
Then $R^*R=RR^*=I_6$ and
$((I-P)RS)^*((I-P)RS)=Q$.
\end{proposition}
\begin{proof}
The positive square roots are commuting functions of $Q$.
Block multiplication cancels the off-diagonal terms and gives
$(I_3-Q)+Q=I_3$ on each diagonal. The lower block of $RS$
is $\sqrt Q$.
\end{proof}

This is a consistency witness, not a recovery of the native pacer.
It also proves why closure, positivity and unitarity alone do not
predict a hierarchy or Koide: all positive contraction spectra remain
available. Any nonzero positive candidate can be scaled into this
range, but that rescaling does not derive a physical mass unit.

\paragraph{Concrete next gate.}
Generation states and an operator-valued coupling now have explicit
conditional formulas: $S=PT(T^*PT)^{-1/2}$ and
$D_\beta=(I-P)\widetilde R_\beta P$, with the computable
three-dimensional Gram matrix $Q_\beta=(D_\beta S)^*D_\beta S$.
What is not supplied is a native evaluated $P$, the action of the
pacer on these seeds, a physical mass-matrix law and any required
relative channel. The cyclic root form explains how a Brannen-type
spectrum can be represented; its Koide radius and phase are still
independent obligations. No lepton-mass fit, native activation,
OP2 closure or parameter-free spectrum is asserted.
